% FEATHER — short conference/introductory paper (Review-I style).
% Purpose: introduce the problem, the motivation, and the direction of the
% project WITHOUT disclosing the method's construction, theory, or results.
% The complete manuscript lives in ../paper/ and stays untouched.
%
% Springer Nature SNmult format, same as the full paper. Compile with
% pdfLaTeX + BibTeX: main.tex + references.bib + SNmult.cls + spmpsci.bst.

\documentclass[graybox]{SNmult}

\usepackage{type1cm}
\usepackage{makeidx}
\usepackage{graphicx}
\usepackage{multicol}
\usepackage[bottom]{footmisc}
\usepackage{newtxtext}
\usepackage[varvw]{newtxmath}

\makeindex

\newcommand{\todo}[1]{\textbf{[TODO: #1]}}

\begin{document}

\title*{Toward Label-Free Detection of Harmful Data Drift in Deployed Classifiers}
\titlerunning{Toward Label-Free Detection of Harmful Drift}
\author{\todo{Name of First Author} and \todo{Name of Second Author}}
\authorrunning{\todo{F. Author and S. Author}}
\institute{\todo{Name of First Author} \at \todo{Name, Address of Institute}, \email{\todo{name@email.address}}
\and \todo{Name of Second Author} \at \todo{Name, Address of Institute}, \email{\todo{name@email.address}}}

\maketitle

\abstract*{Machine learning models degrade after deployment because the data
they receive keeps changing, and the labels that would reveal the damage
arrive late or never. Existing label-free monitors either alarm on every
distribution change, harmful or not, or infer health from the model's own
outputs --- confidence, entropy, uncertainty --- signals that are themselves
unreliable under the very shifts they are meant to detect. This short paper
introduces FEATHER, an ongoing project that approaches the problem from the
geometry of the deployed model itself: rather than asking whether the data
changed, it asks whether the change occurred somewhere the model's behavior
can respond to at all. We motivate the problem, survey the gap in current
practice, state our objectives and evaluation plan on standard drift
benchmarks, and outline the work ahead.
\keywords{Concept drift $\cdot$ Model monitoring $\cdot$ Dataset shift
$\cdot$ Label-free evaluation}}

\abstract{Machine learning models degrade after deployment because the data
they receive keeps changing, and the labels that would reveal the damage
arrive late or never. Existing label-free monitors either alarm on every
distribution change, harmful or not, or infer health from the model's own
outputs --- confidence, entropy, uncertainty --- signals that are themselves
unreliable under the very shifts they are meant to detect. This short paper
introduces FEATHER, an ongoing project that approaches the problem from the
geometry of the deployed model itself: rather than asking whether the data
changed, it asks whether the change occurred somewhere the model's behavior
can respond to at all. We motivate the problem, survey the gap in current
practice, state our objectives and evaluation plan on standard drift
benchmarks, and outline the work ahead.
\keywords{Concept drift $\cdot$ Model monitoring $\cdot$ Dataset shift
$\cdot$ Label-free evaluation}}

\section{Introduction}
\label{sec:intro}

A classifier that was accurate on the day it shipped does not stay accurate
for free. The data it sees in production drifts --- new user behavior, new
sensors, new lighting, new vocabulary --- and the drift is only sometimes
harmless. Practitioners have known this for a long time; Sculley et al.\
counted changing input distributions among the chief sources of technical
debt in real machine learning systems a decade ago \cite{sculley2015debt},
and the problem has since acquired its own literature under the name
\emph{concept drift} \cite{gama2014survey, lu2018learning}.

What makes the problem stubborn is not detecting change. Statistics has
offered change detectors since Page's 1954 CUSUM chart
\cite{page1954cusum}, and the stream-mining community has refined them for
decades \cite{gama2004learning, bifet2007adwin}. The stubborn part is a
mismatch of questions. The operator of a deployed model does not actually
want to know \emph{did the data change?} --- in any real stream the honest
answer is ``constantly.'' The question that matters is \emph{is the model
still working?}, and the labels that would answer it directly arrive late,
sparsely, or never. Every practical monitor is therefore a proxy, and the
choice of proxy is where current practice, in our view, has an unexamined
weakness.

\section{The Gap in Current Practice}
\label{sec:gap}

Label-free monitors in use today fall into two families, and each fails the
operator in its own way.

The first family watches the \emph{inputs}. Two-sample tests such as MMD
\cite{gretton2012kernel}, windowed Kolmogorov--Smirnov statistics
\cite{raab2020kswin}, and adaptive-window detectors \cite{bifet2007adwin}
raise an alarm whenever the arriving distribution departs from the
reference. These tests are sensitive but indiscriminate: they cannot tell a
shift the model absorbs without losing a point of accuracy from one that
quietly halves it. Rabanser et al.\ demonstrated this empirically ---
detection strength and actual damage correlate poorly --- and explicitly
left the characterization of \emph{harmful} shift as an open problem
\cite{rabanser2019failing}. An operator flooded with alarms about harmless
change eventually stops reading them.

The second family watches the \emph{outputs}. Softmax confidence
\cite{hendrycks2017baseline}, entropy, agreement statistics
\cite{guillory2021doc}, learned confidence thresholds \cite{garg2022atc},
prediction-uncertainty indices \cite{pudd2025}, and error-proxy models with
sequential tests \cite{amoukou2024sequential} all estimate the model's
health from the model's own behavior. This is the right question at last,
but the witness is compromised: deep networks are systematically
miscalibrated \cite{guo2017calibration}, and their calibration degrades
further precisely under the distribution shift the monitor is supposed to
catch \cite{ovadia2019trust}. More fundamentally, any statistic computed
from a model's outputs is, by construction, silent about whatever the
outputs themselves cannot express. Whether that silence is a negligible
technicality or a large and structured blind region is, we believe, a
question that has not been asked carefully --- and it is the question our
project starts from.

\section{The FEATHER Project}
\label{sec:project}

FEATHER (Fisher-Eigenvalue Adaptive Thresholding for Error-Robust drift
detection) is our ongoing attempt to build a monitor around the deployed
model's own geometry rather than around its data or its outputs. The
guiding intuition is simple to state. A trained classifier does not use its
internal representation uniformly: some directions in feature space carry
the information its decisions depend on, and others provably do not. The
Fisher information of the model --- a classical object from information
geometry \cite{amari1998natural} that already plays a central role in
optimization \cite{martens2015kfac}, generalization analysis
\cite{karakida2019universal}, and continual learning
\cite{kirkpatrick2017overcoming} --- gives a principled way to make that
distinction quantitative. A monitor built on this distinction could, in
principle, separate the drift a model can be expected to tolerate from the
drift it cannot even see, which is exactly the separation the two existing
families of monitors lack. Related geometric ideas have recently begun to
appear in the monitoring literature \cite{zhang2023score, pandey2025zdp},
which we read as evidence the direction is timely; to our knowledge, no
existing method connects this geometry to the harmful-versus-benign
distinction that operators actually need.

We deliberately keep the construction at this level of description here;
the precise formulation, its theoretical guarantees, and our experimental
findings are the subject of the full manuscript in preparation.

\runinhead{Objectives.} The project has four concrete goals:
\begin{enumerate}
\item{a formal characterization of what output-based monitors can and
  cannot detect for standard classifier architectures;}
\item{a streaming monitor that exploits this characterization, cheap enough
  to run beside production inference on commodity hardware;}
\item{a principled calibration procedure, so that the false-alarm rate is a
  parameter the operator chooses rather than a constant someone tunes; and}
\item{an evaluation protocol in which the harmfulness of a shift is
  \emph{measured} from held-out labels rather than assumed from the name of
  the corruption applied.}
\end{enumerate}

\runinhead{Planned Evaluation.} We will evaluate on the standard public
benchmarks for controlled distribution shift --- rotated digit streams and
the common-corruptions suite CIFAR-10-C \cite{hendrycks2019benchmarking}
--- against representatives of both existing families: classical stream
detectors as implemented in River \cite{montiel2021river}, and
output-based baselines from the confidence-tracking line
\cite{hendrycks2017baseline, garg2022atc}. Alongside detection quality we
will report false-alarm behavior on shifts measured to be benign, which the
literature tends to underweight, and the computational cost of monitoring
relative to the cost of inference itself.

\runinhead{Feasibility.} The project is deliberately scoped to be
achievable within a semester by a two-person team. It requires no data
collection (all benchmarks are public), no large-scale training (the models
involved train in minutes on a single workstation GPU), and the monitoring
computations we envision are orders of magnitude cheaper than the inference
they observe. A synthetic testbed with analytically known answers, which we
have already built, lets us verify each component against ground truth
before any deep network enters the picture.

\section{Future Work}
\label{sec:future}

The immediate roadmap is: complete the theoretical characterization; finish
the monitor implementation and its calibration procedure; run the benchmark
evaluation described above; and write up the results. Beyond the semester,
two directions look promising. First, current corruption benchmarks mostly
produce shifts that output-based monitors can also see; constructing
benchmarks for the \emph{silent} regime --- realistic drift that leaves
model outputs untouched, in the spirit of the naturally occurring shifts
cataloged by WILDS \cite{koh2021wilds} --- would give the community a
sharper instrument than corrupted test sets. Second, the same geometric
viewpoint should extend beyond the classifier head toward deeper layers and
other architectures, including the representational-drift questions now
being asked of large language models \cite{pandey2025zdp}.

\section{Conclusion}
\label{sec:conclusion}

Monitoring a deployed model without labels currently means choosing between
alarms about everything and trust in a witness with a known blind spot. We
have argued that the missing ingredient is the geometry of the deployed
model itself, and introduced FEATHER, a project that develops this
observation into a formal characterization, a practical streaming monitor,
and an evaluation that takes the difference between harmful and benign
drift seriously. The full method, theory, and experimental results will
appear in the complete manuscript.

\begin{acknowledgement}
\todo{Acknowledgements: project guide, institution, and compute support.}
\end{acknowledgement}

\ethics{Competing Interests}{The authors have no conflicts of interest to
declare that are relevant to the content of this chapter.}

\bibliographystyle{spmpsci}
\bibliography{references}

\end{document}
